We introduced \textsc{RECTOR}, a framework that elevates trajectory selection from a simple probability check to a rigorous, rule-governed decision process.  The discussion below first recapitulates the mechanism, then summarizes the empirical evidence, and finally consolidates the paper's principal limitations and deployment qualifications.

%---------- Mechanism recap ----------
\textbf{How RECTOR works.}
Given a set of $K$ candidate ego trajectories produced by a Transformer--CVAE generator, RECTOR applies a three-stage selection pipeline (Fig.~\ref{fig:rector_architecture}).  First, a reference applicability mechanism predicts which of 28~cataloged traffic rules are relevant to the current scene.  In the implemented model this mechanism is a learned head (3.33\,M parameters; 37.7\% of the 8.83\,M-parameter model), pretrained in a dedicated first stage with focal binary cross-entropy ($\gamma{=}2.0$) and then fine-tuned end to end, but the ablations show that a parameter-free always-on policy is preferable for Safety and Legal tiers on this benchmark.  Second, 24~\emph{differentiable proxy functions} score each candidate against the active rules (Table~\ref{tab:proxy_correlation_detailed}; mean pairwise candidate-ranking concordance\,=\,0.948 across all rules, 0.875 for the 10 rules with measurable variance).  Third, a \emph{tiered lexicographic scorer}---with zero trainable parameters---selects the trajectory with the smallest violation vector under the strict priority order \textbf{Safety\,$\succ$\,Legal\,$\succ$\,Road\,$\succ$\,Comfort}, pruning candidates tier by tier and breaking any remaining exact ties by the lowest residual aggregate tier score, with candidate index used only as a final exact-tie fallback (Algorithm~2).  Because the lexicographic ordering is deterministic and parameter-free, higher-priority tiers are \emph{never} traded for lower-tier gains by construction.  An independently tuned weighted-sum baseline (125-vector grid search on the canonical checkpoint; Table~\ref{tab:ws_independent}) confirms, however, that all tested weight configurations converge to the same binary compliance outcomes as RECTOR on 43{,}219 scenarios.  On the present benchmark, the lexicographic machinery is therefore supported as a structural guarantee and audit mechanism, not as an empirically superior compliance rule relative to weighted-sum.

%---------- Evidence ----------
\textbf{Evidence of effectiveness.}
Table~\ref{tab:selection_strategy} and Fig.~\ref{fig:selection_comparison} present the central empirical result: on 12{,}800 WOMD \texttt{validation\_interactive} scenarios, compliance-aware selection raises reported selected-mode compliance from 76.14\% to 84.97\% (equivalently, total violations fall by 8.83~pp, 23.86\%\,$\to$\,15.03\%) and Safety-tier violations by 8.72~pp (21.87\%\,$\to$\,13.15\%) relative to confidence-only selection, while all strategies operate on the \emph{same} candidate set.  Because WOMD is already largely rule-abiding, this is best read as better selection among near-feasible candidates, with the clearest headroom concentrated in the Safety tail, rather than as broad correction of pervasive rule breaking.  Both RECTOR and the uniformly weighted-sum baseline achieve identical compliance outcomes (Table~\ref{tab:selection_strategy}); they differ in selADE (RECTOR: 2.143\,m, weighted-sum: 2.043\,m) because RECTOR's confidence tiebreaker selects different trajectories in tied-compliance cases.

The per-tier breakdown (Table~\ref{tab:per_tier_breakdown}) shows that Safety provides the dominant mode discrimination, while Legal and Road violations are 0.0\% under predicted applicability.  That concentration is the point: most of the measurable gain comes from filtering the residual high-risk cases, not from fixing widespread rule violations.  Compliance-aware selection simultaneously \emph{improves} geometric accuracy: RECTOR lexicographic selADE decreases from 2.208\,m (confidence-only) to 2.143\,m, and the weighted-sum baseline achieves 2.043\,m, because collision-prone, geometrically extreme modes are filtered by Safety-tier enforcement.  The median selADE is 1.93\,m (Table~\ref{tab:error_percentiles}), with the mean driven upward by complex, regulation-heavy scenes.

Ablation studies isolate the selection-time contribution.  The selector ablation (Table~\ref{tab:selector_ablation_protB}) confirms that the tiered scorer is the dominant compliance driver: removing it increases Safety-tier violations by 6.07 percentage points \emph{with zero change in candidate-set accuracy}, confirming that compliance gains stem from the selection rule rather than from changes to trajectory generation.  RECTOR's selection-time approach preserves full multimodality: 6.0 effective modes versus 1.0--2.3 for constrained-generation methods (Table~\ref{tab:mode_diversity}), demonstrating that compliance need not sacrifice the diversity required for downstream planning under uncertainty.

A 55-scenario Waymax log-replay validation achieves zero overlap steps across 3{,}338.5\,s of simulated driving (Section~\ref{subsec:closed_loop}), confirming that the integration pipeline executes the selected trajectories end to end.

From an engineering standpoint, the framework adds minimal overhead: mean inference latency is 7.3\,ms per sample (136.5~samples/s), with the rule-evaluation and selection stage contributing only 1.4\,ms---19\% of the total pipeline (Table~\ref{tab:efficiency}).  This is faster than Wayformer (85~samples/s) and comparable to DenseTNT (120~samples/s), confirming real-time viability.  The verification pipeline (Table~\ref{tab:semantic_invariants})---semantic invariants (99.2--100\% pass), numerical stability (0.0\% NaN/Inf), and proxy--evaluator pairwise concordance (0.948 all rules; 0.875 for rules with measurable variance)---ensures that the reported improvements are not artifacts of evaluator instability, and every selection decision produces a fully auditable trace: active rules, per-tier scores, and the first tier that differentiated candidates.

%---------- Limitations ----------
\textbf{Principal limitations and qualifications.}
The results above should be read with six consolidated qualifications.  First, because the WOMD benchmark is already largely rule-abiding, the headline 13.15\% Safety-tier violation rate leaves limited absolute headroom for improvement; under oracle applicability, Safety violations rise to 19.94\% and total violations to 32.11\%, so the predicted-applicability numbers are optimistic but still describe selection among a mostly compliant candidate pool.  Second, the applicability-head study is best interpreted as a negative result with immediate practical implications: the 3.33\,M-parameter learned head is strictly dominated by a parameter-free always-on policy for Safety and Legal tiers, so the learned component should not be treated as the main practical driver of the system on this benchmark.  Third, the weighted-sum baseline matches RECTOR on the present evaluation set and achieves lower selADE, so the paper's demonstrated empirical gain is the improvement over confidence-only selection; RECTOR's additional contribution here is the structural enforcement of the tier ordering and the accompanying audit trace, not a measured compliance advantage over weighted-sum.  Fourth, the current benchmark does not expose a regime where weight misspecification changes binary compliance outcomes; demonstrating that use case requires additional evaluation under distribution shift, adversarial weight perturbations, or scenarios with stronger cross-tier tension.  Fifth, the current end-to-end study uses a single Transformer--CVAE generator, so generator-agnostic composability is an architectural property rather than a multi-generator empirical claim.  Sixth, the evaluation is open-loop on WOMD with a 5-second horizon and U.S.-centric rule definitions; the 55-scenario Waymax study validates the simulator bridge with a mock log-replay generator rather than a full reactive closed-loop comparison, and the current stack still has incomplete proxy coverage for four rules plus measurable sensitivity to map and signal errors.  Together, these limitations define the current boundary of the contribution: RECTOR demonstrates that explicit priority-ordered selection can materially improve compliance on a fixed candidate set, but additional work on applicability policy design, generator quality, closed-loop interaction, proxy coverage, robustness, and distribution-shift evaluation is still required before production use.
